\scanpage{11}
% L1–3
\deleted{\unsure{ce qui|où pour}}\inserted{où} pour tout $i\in\unsure{\tau|T}\amalg\tau'\deleted{\unsure{P}}=P(V)$, $P_i$ est un $D_i$-torseur, défini par
\[
P_i=P/D_{\bar i}=P\wedge_V(D_i\simeq V/D_{\bar i})
\tag*{(6)}
\]
% L4–9
où $i\mapsto\bar i$ est l'\deleted{\unsure{élément}}\inserted{involution} de $P(V)$ qui \deleted{est} correspond à sa partition de type $(2,2)$ $\{\tau,\tau'\}$, qui dans $\tau$ et dans $\tau'$ se réduit à la transposition dans ces ensembles. Ainsi, l'ensemble $P$, de cardinal $9$, est muni, via
\[
P\simeq\prod_{i\in\tau}P_i
\tag*{$(*)$}
\]
% L10–18
\deleted{de}d'\deleted{deux}\inserted{une \deleted{\unsure{autre}}paire de} partitions de type $(3,3,3)$, \unsure{indexée|indexées} par $\tau$, ces partitions étant ``d'équerre'', en ce sens que l'application canonique \inserted{\unsure{ces|cs|de}} de $P$ sur le produit des deux ens. quotients correspondants est bijective, ce qui signifie aussi que pour tte classe de l'une et toute classe de l'autre, l'intersection a exactement un élément (il suffit de supposer qu'elle est tjrs non vide; ou qu'elle a tjrs au plus un élément\ldots). Et ce même ensemble $P$ est muni d'une deuxième paire (distincte de la première) de partitions de type $(3,3)$, \inserted{``en équerre'', \unsure{indexée|indexées} par $\tau'$,} correspondant à la décomposition
\[
P\simeq\prod_{i\in\tau'}P_i
\tag*{$(*')$}
\]
% archive sheet number: 10

